\chapter{75}

\section{Moosfluh, Bettmeralp (CH), Dienstag, 4.
September}



Die vielen Aufgaben liessen die Stunden erstaunlich schnell vergehen.
Selbst die gemeinsamen Essenszeiten, kurz gehalten, boten kaum Raum,
Persönliches zu erzählen. Louis sollte es recht sein. Und auch heute
Nacht würde er es zum dritten Mal schaffen, Ava im gemeinsamen
Schlafraum nicht näher zu begegnen. Er wollte später schlafen und auch
früher aufstehen als sie. Nur keine neuen, zusätzlichen Verwicklungen.
Keine Fragen. Das Auftauchen von Bastianos Gesicht, in den unmöglichsten
Momenten, war ihm Schrecken genug.
Während des Aufstiegs auf die Alp würde er Zeit finden, Valentina
anzurufen. Was wollte Bastiano bei ihr? Warum hatten die Kontakt? Was
wusste sie von dem?

\sectionbreak

Der junge Tag hatte etwas Knackiges, Unverbrauchtes, und von sich aus hätte er
niemandem etwas angetan. Louis war früh dran. Ihm schien, als läge die Wiese noch im Tiefschlaf 
und er, als Aufsteigender mittendrin. Andächtig setzte er Schuh vor Schuh. Gleichmässig und ruhig. 
Wenn es nichts vorher und nichts nachher gäbe, wäre gerade alles in Ordnung. Bis zum Durchbruch der ersten, feinen Sonnenstrahlen trug er den Gedanken bei sich. Mit der Tagessicht kam unausweichlich die Realität zurück. Giorgia tauchte auf. Wie sie in ihrem Spitalbett lag und nach Erinnerungen suchte. Und keine fand. 
Der Aufstieg war steil. Louis atmete allmählich schwerer. In gut
sichtbarer Entfernung zeichnete sich ein kleines Plateau ab. Eine erste
Pause und das Gespräch mit Valentina standen an. Wenn sie denn
erreichbar war. Angestrengt häufte er mögliche Varianten an, den Anruf
zu beginnen. Selbst wenn es noch angenehm kühl war, fing er an zu
schwitzen. Und wieder begannen seine Gedanken zu rasen. Ihm war, als
stelle sich Bastianos Standbild wieder vor ihm auf. Unverrückbar
hartnäckig. Es war lästig. Was hatte der mit Valentina zu tun? 
Er wählte ihre Nummer.

«Louis, hallo,» Valentina hörte sich gelöst an.

«Valentina, hi, stör ich dich?»

«Nein, passt gut. Ich habe gerade an dich gedacht.»

Ihr freudiger Unterton öffnete Louis augenblicklich eine Tür.

«Schön. Ehm, nach unserem Gespräch gestern habe ich über einiges
nachgedacht, und,» Louis hüstelte, «und, ich war erstaunt, dass Bastiano
zu dir kam?»

«Ja, du, ich auch,» sie machte eine Pause, «jetzt sehe ich, du riefst
mich gestern nochmals an?»

«Ja, das heisst, nein, ich habe aus Versehen die Rückruftaste gedrückt,
alles okay.»

«Aha, gut, also, Basti wirkte auf mich eingeschnappt. Ich liess ihn
rein, ging kurz zur Toilette, um mich zu sammeln, und als ich zurückkam,
sass der auf meinem Sofa, stand aber sofort auf und starrte mich fragend
an.»

«Was wollte er denn von dir?»

«Das wusste ich erst auch nicht. Zuerst sagte er ausgesprochen wütend,
ich wisse bestimmt, wo du seist, ob ich dich decke. Ich fragte, wie er
denn darauf komme. Er reagierte nicht auf meine Frage, sondern redete
von Giorgia, von Einvernahmen auf der \emph{Questur}, davon, dass die Polizei
ihn verdächtige, mit Giorgias Sturz etwas zu tun zu haben, wieder von
irgendwelchen Spuren, davon, dass sein Vater ihn rausholen musste, und
dass er nichts getan habe, und dass die ihn gefragt hätten, ob er wisse,
wo du seist, ob er ein Freund sei von dir,» Valentina Stimme brach, «und
dann kam er wieder damit, warum er wirklich hier war. Er sagte, ich sei
ja gut mit dir befreundet und er sei sicher, dass ich wisse, wo du bist
und, er habe gehört, du seist vor Giorgias Sturz in Richtung Felsen ins
Wäldchen verschwunden.»

Valentina war nach und nach leiser geworden. Louis setzte sich
straff hoch.

«Ach ja? Und von wem will der das gehört haben?»

Ihm blieb die Luft weg.

«Das wollte er mir nicht verraten. Ich sagte dann nochmal, wir hätten
gerade Funkstille, und dass ich mich auch frage, wo du bist,» sie atmete
hörbar aus, «stimmt ja eigentlich. Keine Sorge, ich habe komplett dicht
gehalten. Dann stand er wie aus dem Nichts auf, bedrückt und genervt und
ging wieder. Louis, ich glaube, der war komplett benebelt.»

«Sieht ganz danach aus. Ich danke dir, Valentina. Und wie ging es dir
danach?»

«Ich war durcheinander, ich fühlte mich überfallen. Es war unwirklich,
hat mich komplett fertig gemacht. Ich habe keine Ahnung, was Basti damit zu tun haben
soll.»

«Ich auch nicht.»

Die Ratlosigkeit schwappte nervös zwischen ihren Handys hin und her.
